\pagestyle{empty}
\tight
\begin{tblr}{width=\textwidth,colspec={|Q[c,m,1.9cm]|X[l,m]|},hlines,vlines,hspan=minimal,row{1}={bg=headblue},rowsep=1pt,colsep=3pt}
\SetCell{c,m}\hfil\logo\hfil & \textbf{POLITEKNIK NEGERI MALANG}\par\textbf{JURUSAN TEKNOLOGI INFORMASI}\par\textbf{PROGRAM STUDI: D4 TEKNIK INFORMATIKA} \\
\SetCell[c=2]{c} \textbf{RENCANA PEMBELAJARAN SEMESTER (RPS) --- DRAF KONSEP C: SEIMBANG} & \\
\end{tblr}
\begin{longtblr}{width=\textwidth,colspec={|Q[l,m,1.9cm]|X[0.85,l,m]|X[1.45,l,m]|X[0.75,l,m]|X[1.15,l,m]|},hlines,vlines,hspan=minimal,rowsep=1pt,colsep=2pt,row{1,3}={bg=headgray,font=\bfseries}}
MATA KULIAH & KODE & BOBOT (SKS)/jam & SEMESTER & TGL. PENYUSUNAN \\
Kuliah Kerja Nyata Tematik & RTI257002 (sementara)\footnotemark[1] & 2 SKS / 4 jam/minggu & 6/7 (sementara)\footnotemark[1] & 13 Juli 2026 \\
OTORISASI & Dosen Pengembang RPS & Koordinator RMK & \SetCell[c=2]{l} Ka PRODI &  \\
 & Tim Pengajar Program Studi D4 TI &  & \SetCell[c=2]{l}{\vspace{8mm}\par Dr. Ely Setyo Astuti, S.T., M.T.} &  \\
\SetCell[r=14]{l} Capaian Pembelajaran (CP) & \SetCell[c=4]{l,bg=headgray,font=\bfseries} Capaian Pembelajaran Lulusan Program Studi (CPL-Prodi) &  &  &  \\
 & CPL01 & \SetCell[c=3]{l} Menginternalisasi ketakwaan kepada Tuhan Yang Maha Esa, menaati hukum, disiplin, dan menunjukkan profesionalisme melalui etika profesi, pembelajaran sepanjang hayat, serta respons terhadap isu sosial dan teknologi. &  &  \\
 & CPL03 & \SetCell[c=3]{l} Mampu mengelola tim, manajemen diri, berkomunikasi secara lisan maupun tertulis dengan baik dan mampu melakukan presentasi. &  &  \\
 & CPL06 & \SetCell[c=3]{l} Mampu merancang, mengimplementasikan, menguji, melakukan deployment, dan memelihara, serta menjamin mutu perangkat lunak yang menjawab permasalahan dan memenuhi kebutuhan stakeholder. &  &  \\
 & CPL08 & \SetCell[c=3]{l} Mampu mengelola proyek teknologi informasi secara profesional dengan mengoptimalkan sumber daya yang tersedia. &  &  \\
 & \SetCell[c=4]{l,bg=headgray,font=\bfseries} Capaian Pembelajaran mata kuliah (CPL-MK) &  &  &  \\
 & CPMK0102 & \SetCell[c=3]{l} Mampu menginternalisasi nilai ketakwaan kepada Tuhan YME, Pancasila, kewarganegaraan, dan menunjukkan pembelajaran sepanjang hayat melalui disiplin diri, tanggung jawab sosial, dan adaptasi karir profesional. &  &  \\
 & CPMK0301 & \SetCell[c=3]{l} Mampu mengelola diri secara profesional, berkomunikasi secara lisan maupun tertulis, dan melakukan presentasi secara efektif dalam konteks kerja. &  &  \\
 & CPMK0803 & \SetCell[c=3]{l} Mampu melaksanakan, memonitor, dan mengendalikan proyek serta mengelola dinamika tim secara efektif. &  &  \\
 & CPMK0607 & \SetCell[c=3]{l} Mampu membangun solusi sistem informasi terintegrasi yang kompleks dan menjawab kebutuhan \textit{stakeholder}. &  &  \\
 & \SetCell[c=4]{l,bg=headgray,font=\bfseries} Kemampuan akhir tiap tahapan belajar (Sub-CPMK) &  &  &  \\
 & SCPMK0102-06001 & \SetCell[c=3]{l} Mahasiswa mampu mengidentifikasi kebutuhan sosial mitra/komunitas dan berkontribusi nyata dalam program pemberdayaan. &  &  \\
 & SCPMK0301-06002 & \SetCell[c=3]{l} Mahasiswa mampu menggali kebutuhan mitra secara partisipatif dan mengomunikasikan hasil program kerja serta luaran KKN. &  &  \\
 & SCPMK0803-06003 & \SetCell[c=3]{l} Mahasiswa mampu menyusun dan melaksanakan rencana program kerja KKN sesuai jadwal yang disepakati bersama mitra. &  &  \\
 & SCPMK0607-06004 & \SetCell[c=3]{l} Mahasiswa mampu merancang dan mewujudkan luaran KKN (solusi TI dan/atau kegiatan pemberdayaan) yang sesuai kebutuhan mitra. &  &  \\
\SetCell[r=2]{l} Penilaian & \SetCell[c=4]{l} *Nilai CPMK didapat dari agregasi nilai SUB-CPMK yang dibebankan &  &  &  \\
 & \SetCell[c=4]{l}{\tiny\begin{tblr}{width=\linewidth,colspec={|Q[c,m,0.1\linewidth]|Q[c,m,0.16\linewidth]|X[l,m]|X[l,m]|X[l,m]|X[l,m]|Q[c,m,0.08\linewidth]|},hlines,vlines,hspan=minimal,rowsep=1pt,colsep=0.7pt,row{1,2}={bg=headgray,font=\bfseries}}
Kode CPMK & Kode SCPMK & \SetCell[c=4]{c} Bobot per Bentuk Penilaian & & & & Total Bobot \\
 & & Proposal & Pelaksanaan \& Keterlibatan & Luaran KKN & Laporan \& Seminar & \\
CPMK0102 & SCPMK0102-06001 & 5\%: etika \& kelayakan program & 15\%: keterlibatan \& kontribusi pemberdayaan & & & 20\% \\
CPMK0301 & SCPMK0301-06002 & 10\%: komunikasi kebutuhan dengan mitra & & & 10\%: komunikasi hasil kepada mitra dan penguji & 20\% \\
CPMK0803 & SCPMK0803-06003 & 5\%: perencanaan \& penjadwalan program kerja & 15\%: pelaksanaan \& pengelolaan program & & 10\%: dokumentasi pelaksanaan \& evaluasi & 30\% \\
CPMK0607 & SCPMK0607-06004 & & & 20\%: rancang bangun luaran KKN & 10\%: defense luaran KKN & 30\% \\
\SetCell[c=6]{r}\textbf{Total Per Penilaian} & & 20\% & 30\% & 20\% & 30\% & \textbf{100\%} \\
\end{tblr}} &  &  &  \\
Diskripsi Singkat Mata Kuliah & \SetCell[c=4]{l} Kuliah Kerja Nyata Tematik (KKN Tematik) adalah mata kuliah wajib 2 SKS penunjang IKU kegiatan mahasiswa di luar Prodi (dasar: Nota Dinas Wakil Direktur I No. 11/WADIR I/KM/2026 dan Permendiktisaintek No. 39 Tahun 2025) yang menempatkan mahasiswa pada mitra/komunitas untuk melaksanakan program kerja pemberdayaan sekaligus menghasilkan luaran (solusi TI dan/atau kegiatan) yang menjawab kebutuhan mitra, disertai laporan dan seminar hasil. \\
Bahan Kajian / Materi Pembelajaran & \SetCell[c=4]{l}\begin{enumerate}\item Pembekalan KKN dan pemetaan kebutuhan mitra/komunitas\item Analisis kebutuhan dan perumusan luaran KKN\item Pelaksanaan program kerja: kegiatan pemberdayaan dan pengembangan luaran KKN\item Pengelolaan tim, komunikasi, dan evaluasi dampak program\item Penyusunan laporan KKN dan seminar hasil\end{enumerate} \\
\SetCell[r=2]{l} Pustaka & \SetCell[c=4]{l} \textbf{Utama:}\begin{enumerate}\item Buku Pedoman KKN Tematik Politeknik Negeri Malang.\item Panduan MBKM Kemdikbud 2020.\item Pedoman Penulisan Laporan KKN Polinema.\end{enumerate} \\
 & \SetCell[c=4]{l} \textbf{Pendukung:} Materi LMS, dokumen profil mitra/komunitas, dan bahan bimbingan dosen pembimbing lapangan (DPL). \\
Media Pembelajaran & \SetCell[c=2]{l} \textbf{Software:} Aplikasi logbook digital, LMS, perangkat bantu pengembangan luaran KKN. &  & \SetCell[c=2]{l} \textbf{Hardware:} Komputer/laptop, perangkat lapangan mitra, perangkat presentasi. &  \\
Nama Dosen Pengampu & \SetCell[c=4]{l} Tim Pengajar Program Studi D4 TI \\
Matakuliah Syarat & \SetCell[c=4]{l} Lulus minimal 90 SKS \\
\end{longtblr}
\begin{longtblr}{width=\textwidth,colspec={|Q[c,m,0.06\textwidth]|X[1.35,l,m]|X[1.08,l,m]|X[1.16,l,m]|X[0.7,l,m]|X[1.24,l,m]|X[1.06,l,m]|X[0.98,l,m]|Q[c,m,0.05\textwidth]|},hlines,vlines,hspan=minimal,rowhead=2,row{1,2}={bg=headgreen,font=\bfseries},rowsep=1pt,colsep=1.5pt}
Mgg.\par Ke & Kemampuan Akhir Yang Direncanakan (Sub-CP-MK) & Bahan kajian (Materi Pembelajaran) & Bentuk dan Metode Pembelajaran & Estimasi Waktu & Pengalaman Belajar Mahasiswa & Kriteria \& Bentuk Penilaian & Indikator Penilaian & Bobot Penilaian (\%) \\
(1) & (2) & (3) & (4) & (5) & (6) & (7) & (8) & (9) \\
1 & SCPMK0102-06001, SCPMK0803-06003 & Pembekalan KKN: orientasi, etika kerja, dan kontrak belajar. & Modalitas: Blended Learning. Bentuk: Luring/Daring. Metode: Case Method, ceramah interaktif, dan diskusi. & 1 x 2 x 50' tatap muka; persiapan mandiri. & Mahasiswa mengikuti pembekalan KKN, memahami etika kerja dan tanggung jawab sosial, serta menandatangani kontrak belajar. & Bentuk penilaian: Kehadiran dan partisipasi pembekalan. Mengacu rubrik RTM. & Ketepatan pemahaman prosedur; kesiapan etika kerja. & 2 \\
2 & SCPMK0301-06002, SCPMK0803-06003 & Observasi lapangan: pemetaan kebutuhan mitra secara teknis dan sosial. & Modalitas: Praktik Lapangan dan Bimbingan. Metode: observasi lapangan, wawancara mitra, dan bimbingan dosen pembimbing lapangan (DPL). & 1 x 2 x 50' observasi lapangan; kerja mandiri. & Mahasiswa melakukan observasi ke lokasi mitra/komunitas dan memetakan kebutuhan baik dari sisi sosial maupun potensi solusi TI. & Bentuk penilaian: Catatan observasi lapangan. Mengacu rubrik RTM. & Ketepatan identifikasi kebutuhan; kelengkapan data lapangan. & 3 \\
3 & SCPMK0301-06002, SCPMK0607-06004 & Analisis kebutuhan mitra dan perumusan luaran KKN. & Modalitas: Praktik Lapangan dan Bimbingan. Metode: analisis partisipatif dan bimbingan DPL. & 1 x 2 x 50' bimbingan; analisis mandiri. & Mahasiswa menganalisis kebutuhan mitra secara partisipatif dan merumuskan luaran KKN yang akan dihasilkan (solusi TI dan/atau kegiatan pemberdayaan). & Bentuk penilaian: Draft analisis kebutuhan dan rumusan luaran. Mengacu rubrik RTM. & Ketepatan analisis kebutuhan; kesesuaian luaran dengan konteks mitra. & 3 \\
4 & SCPMK0102-06001, SCPMK0803-06003 & Penyusunan proposal program kerja. & Modalitas: Blended Learning. Bentuk: Luring/Daring. Metode: Case Method, penyusunan dokumen, dan bimbingan DPL. & 1 x 2 x 50' bimbingan; penyusunan mandiri. & Mahasiswa menyusun proposal program kerja KKN yang memuat rencana kegiatan, jadwal, dan rancangan luaran KKN bagi mitra. & Bentuk penilaian: Draft proposal program kerja. Mengacu rubrik RTM. & Kelengkapan proposal; kelayakan rencana kerja; ketepatan etika penyusunan. & 5 \\
5 & SCPMK0102-06001, SCPMK0301-06002, SCPMK0803-06003 & Seminar proposal program kerja. & Modalitas: Blended Learning. Bentuk: Luring/Daring. Metode: presentasi dan Case Method. & 1 x 2 x 50' seminar proposal. & Mahasiswa mempresentasikan proposal program kerja KKN di hadapan DPL dan menerima masukan untuk perbaikan. & Bentuk penilaian: Seminar proposal program kerja. Mengacu rubrik RTM. & Kelayakan program kerja; kejelasan rancangan luaran KKN; kualitas presentasi. & 15 \\
6 & SCPMK0803-06003, SCPMK0607-06004 & Pelaksanaan program kerja tahap 1: kegiatan pemberdayaan dan rancangan luaran KKN. & Modalitas: Praktik Lapangan dan Bimbingan. Metode: praktik lapangan, monitoring berkala, dan bimbingan DPL. & 1 x 2 x 50' bimbingan; kerja lapangan mandiri. & Mahasiswa melaksanakan kegiatan pemberdayaan awal di lokasi mitra sekaligus merancang luaran KKN sesuai kebutuhan yang telah dianalisis. & Bentuk penilaian: Logbook dan draft rancangan luaran. Mengacu rubrik RTM. & Ketepatan waktu; kesesuaian rancangan; kualitas dokumentasi. & 3 \\
7 & SCPMK0607-06004 & Pelaksanaan program kerja tahap 1 (lanjutan): implementasi awal luaran KKN. & Modalitas: Praktik Lapangan dan Bimbingan. Metode: praktik lapangan, monitoring berkala, dan bimbingan DPL. & 1 x 2 x 50' bimbingan; kerja lapangan mandiri. & Mahasiswa mengimplementasikan bagian awal luaran KKN dan menguji cobakannya bersama mitra. & Bentuk penilaian: Logbook dan progress implementasi. Mengacu rubrik RTM. & Kualitas implementasi; ketepatan waktu; keterlibatan mitra. & 3 \\
8 & SCPMK0102-06001, SCPMK0803-06003 & Monitoring dan evaluasi (monev) tengah bersama mitra dan DPL. & Modalitas: Praktik Lapangan dan Bimbingan. Metode: praktik lapangan, monitoring berkala, dan bimbingan DPL. & Disesuaikan jadwal mitra. & Mahasiswa menjalani monev tengah program KKN bersama mitra dan DPL untuk mengevaluasi kemajuan program kerja dan luaran KKN. & Bentuk penilaian: Formulir monev tengah oleh mitra/DPL. Mengacu rubrik RTM. & Kesesuaian progres dengan rencana; profesionalisme; kualitas kemajuan luaran. & 15 \\
9 & SCPMK0607-06004 & Pelaksanaan program kerja tahap 2: penyempurnaan luaran KKN berdasarkan masukan monev. & Modalitas: Praktik Lapangan dan Bimbingan. Metode: praktik lapangan, monitoring berkala, dan bimbingan DPL. & 1 x 1 x 50' monitoring; kerja mandiri. & Mahasiswa menyempurnakan luaran KKN berdasarkan masukan hasil monev tengah. & Bentuk penilaian: Logbook dan revisi luaran. Mengacu rubrik RTM. & Kualitas revisi; ketepatan waktu. & 3 \\
10 & SCPMK0301-06002, SCPMK0607-06004 & Pelaksanaan program kerja tahap 2 (lanjutan): kegiatan pemberdayaan dan pelatihan penggunaan luaran KKN. & Modalitas: Praktik Lapangan dan Bimbingan. Metode: praktik lapangan, monitoring berkala, dan bimbingan DPL. & 1 x 1 x 50' monitoring; kerja mandiri. & Mahasiswa melaksanakan kegiatan pemberdayaan lanjutan dan melatih mitra menggunakan luaran KKN yang dihasilkan. & Bentuk penilaian: Logbook dan berita acara pelatihan. Mengacu rubrik RTM. & Kualitas pelatihan; ketepatan waktu; keterlibatan mitra. & 3 \\
11 & SCPMK0803-06003 & Pelaksanaan program kerja tahap 3: evaluasi dampak program. & Modalitas: Praktik Lapangan dan Bimbingan. Metode: praktik lapangan, monitoring berkala, dan bimbingan DPL. & 1 x 1 x 50' monitoring; kerja mandiri. & Mahasiswa mengevaluasi dampak program kerja dan luaran KKN terhadap komunitas/mitra. & Bentuk penilaian: Logbook dan catatan evaluasi dampak. Mengacu rubrik RTM. & Ketepatan waktu; kualitas evaluasi dampak. & 3 \\
12 & SCPMK0607-06004 & Konsolidasi luaran dan dokumentasi. & Modalitas: Praktik Lapangan dan Bimbingan. Metode: praktik lapangan, monitoring berkala, dan bimbingan DPL. & 1 x 1 x 50' monitoring; kerja mandiri. & Mahasiswa menyusun dokumentasi luaran KKN (teknis dan/atau kegiatan) sebagai hasil akhir program. & Bentuk penilaian: Dokumentasi luaran KKN. Mengacu rubrik RTM. & Kelengkapan dokumentasi; ketepatan waktu. & 3 \\
13 & SCPMK0102-06001, SCPMK0301-06002 & Bimbingan penyusunan laporan KKN. & Modalitas: Blended Learning. Bentuk: Luring/Daring. Metode: Case Method dan bimbingan DPL. & 1 x 1 x 50' bimbingan; penulisan mandiri. & Mahasiswa mendapatkan bimbingan teknis penulisan laporan KKN mencakup sistematika, konten, dan tata tulis ilmiah. & Bentuk penilaian: Draft laporan dengan catatan bimbingan. Mengacu rubrik RTM. & Ketepatan waktu; kualitas dokumen. & 5 \\
14 & SCPMK0301-06002 & Finalisasi laporan KKN. & Modalitas: Blended Learning. Bentuk: Luring/Daring. Metode: Case Method dan bimbingan DPL. & 1 x 1 x 50' bimbingan; penulisan mandiri. & Mahasiswa menyelesaikan draft akhir laporan KKN mencakup seluruh bab dan lampiran. & Bentuk penilaian: Laporan KKN final. Mengacu rubrik RTM. & Kelengkapan laporan; ketepatan waktu. & 5 \\
15 & SCPMK0102-06001, SCPMK0803-06003 & Persiapan seminar hasil: penyusunan bahan presentasi dan kelengkapan dokumen. & Modalitas: Blended Learning. Bentuk: Luring/Daring. Metode: Case Method dan bimbingan DPL. & 1 x 1 x 50' persiapan; mandiri. & Mahasiswa menyiapkan bahan presentasi seminar hasil dan melengkapi dokumen pendukung program kerja. & Bentuk penilaian: Kelengkapan dokumen seminar. Mengacu rubrik RTM. & Kelengkapan dokumen; kesiapan presentasi. & 5 \\
16 & SCPMK0301-06002, SCPMK0607-06004 & Seminar hasil KKN: presentasi dan defense laporan serta luaran KKN. & Modalitas: Blended Learning. Bentuk: Luring/Daring. Metode: presentasi dan Case Method. & 1 x 2 x 50' seminar hasil. & Mahasiswa mempresentasikan dan mempertahankan laporan KKN beserta luaran KKN yang dihasilkan di hadapan tim penguji. & Bentuk penilaian: Seminar hasil (presentasi dan tanya jawab). Mengacu rubrik RTM. & Kualitas presentasi; ketepatan jawaban; kualitas luaran KKN. & 25 \\
\end{longtblr}
\footnotetext[1]{\textit{Draf konsep untuk perbandingan. Nama mata kuliah dan bobot 2 SKS mengikuti Nota Dinas Wakil Direktur I No. 11/WADIR I/KM/2026. Kode mata kuliah, semester, dan integrasi ke crosswalk kurikulum akan difinalkan pada versi yang dipilih.}}
